\chapter{METODE PENELITIAN}

\section{Desain Penelitian}
Penelitian dilakukan dengan menggunakan metode Design and Development Research (DDR). DDR menurut \textcite{richeyDesignDevelopmentResearch2014c} adalah metode penelitian yang secara sistematis mempelajari, mengembangkan, dan mengevaluasi produk, alat, atau model untuk menghasilkan solusi praktis sekaligus pengetahuan teoritis yang berbasis empiris. Tahapan dalam DDR terdiri dari tiga fase utama, yaitu analisis kebutuhan, perancangan dan pengembangan, serta evaluasi dan pengujian. Gambar \ref{fig:ddr} memberikan visualisasi dari tahapan dalam DDR. Metode ini digunakan sesuai dengan tujuan penelitian yang mengembangkan dan mengevaluasi sistem agentik untuk identifikasi penyakit tanaman sebagai fokus utama.

\begin{figure}[H]
    \centering
    \includegraphics[width=0.8\textwidth]{bab3/ddr.png}
    \caption{Fase dalam penelitian Design and Development Research}
    \label{fig:ddr}
\end{figure}

\section{Prosedur Penelitian}
Prosedur dari penelitian ini dibagi menjadi beberapa tahapan berdasarkan metode penelitian yang digunakan. Gambar \ref{fig:prosedur_penelitian} memberikan visualisasi dari tahapan penelitian yang diatur secara terurut dan sistematis. Penelitian diawali dengan studi literatur mendalam terhadap penelitian-penelitian terkait untuk mengidentifikasi kesenjangan dan tantangan dalam penerapan AI pada pengembangan sistem untuk identifikasi penyakit tanaman. Tahap berikutnya adalah analisis kebutuhan yang mencakup teknik, pendekatan, dan metode yang relevan untuk mendukung pengembangan sistem. Hasil dari analisis kebutuhan ini menjadi dasar dalam integrasi YOLO, VLM, dan perancangan sistem agentik dengan pendekatan agen LLM. Berdasarkan rancangan tersebut, dilakukan perancangan dan pengembangan modul-modul yang membentuk sistem secara keseluruhan. Modul-modul tersebut meliputi modul deteksi objek tanaman, modul pencarian gambar-teks penyakit tanaman, dan modul pencarian informasi.

Modul-modul yang ada akhirnya diintegrasi menjadi sistem agentik yang terpadu. Sistem yang terintegrasi akhirnya dapat diuji dan dievaluasi untuk mengukur performa dan fungsionalitas yang dihasilkan. Sistem juga akan diimplementasikan pada sebuah aplikasi web sebagai antarmuka.

\begin{figure}[H]
    \centering
    \includegraphics[width=0.8\textwidth]{bab3/prosedurpenelitian.png}
    \caption{Tahapan dalam prosedur penelitian}
    \label{fig:prosedur_penelitian}
\end{figure}

\section{Analisis Kebutuhan}
Tabel \ref{tab:perangkat_pendukung} menyajikan perangkat pendukung yang digunakan dalam penelitian ini.

\begin{table}[H]
    \centering
    \caption{Perangkat Pendukung Penelitian}
    \label{tab:perangkat_pendukung}
    \small
    \setstretch{1.0}
    \begin{tabular}{|c|l|l|p{7cm}|}
    \hline
    \multicolumn{4}{|c|}{\textbf{Spesifikasi Laptop}} \\ \hline
    \textbf{No.} & \textbf{Komponen} & \multicolumn{2}{l|}{\textbf{Spesifikasi}} \\ \hline
    1. & Prosesor & \multicolumn{2}{l|}{AMD Ryzen™ 7 7840HS} \\ \hline
    2. & Memori & \multicolumn{2}{l|}{16 GB DDR5-5600} \\ \hline
    3. & Grafik & \multicolumn{2}{l|}{NVIDIA® GeForce RTX™ 4060 8GB} \\ \hline
    \multicolumn{4}{|c|}{\textbf{Perangkat Lunak}} \\ \hline
    \textbf{No.} & \textbf{Nama} & \textbf{Versi} & \textbf{Fungsi} \\ \hline
    1. & WSL2 & - & Subsistem Linux untuk Windows. \\ \hline
    2. & Docker & - & Platform kontainerisasi aplikasi. \\ \hline
    3. & uv & - & Manajer paket dan lingkungan Python. \\ \hline
    4. & Streamlit & - & Framework antarmuka aplikasi web. \\ \hline
    5. & LangChain & - & Framework pengembangan aplikasi LLM. \\ \hline
    6. & LangGraph & - & Library untuk alur kerja agen berbasis graf. \\ \hline
    8. & LangSmith & - & Platform monitoring dan evaluasi LLM. \\ \hline
    9. & AgentEval & - & Framework evaluasi agen cerdas. \\ \hline
    10. & DeepEval & - & Unit testing untuk aplikasi LLM. \\ \hline
    11. & RAGAS & - & Metrik evaluasi sistem RAG. \\ \hline
    12. & Qdrant & - & Database vektor untuk pencarian semantik. \\ \hline
    13. & Google Gen AI & - & API model bahasa Gemini. \\ \hline
    14. & Voyage AI & - & API model embedding dan reranker. \\ \hline
    \end{tabular}
\end{table}
% TODO(AI): Perlu update isi tabelnya nanti.

Sistem agentik dalam penelitian ini dirancang dengan mengadopsi pola \textit{Reasoning and Acting} (ReAct). Implementasi sistem dilakukan menggunakan framework LangChain melalui library LangGraph \parencite{chaseLangChain2022}. Arsitektur sistem yang dikembangkan bersifat multimodal, di mana agen mampu menerima dan memproses input berupa citra tanaman serta instruksi tekstual dari pengguna secara bersamaan. Gambar \ref{fig:arsitektur_sistem_agentik} menyajikan arsitektur sistem agentik beserta spesifikasi input dan output pada setiap alat yang terintegrasi.

%TODO: input gambar nyambung juga harusnya ke owlv2, terus output harusnya bukan cropped images, tapi skor, bounding box, dan gambar teranotasi
\begin{figure}[H]
    \centering
    \includegraphics[width=1\textwidth]{bab3/arsitektursistemagentik.png}
    \caption{Arsitektur sistem agentik dengan kemampuan multimodal}
    \label{fig:arsitektur_sistem_agentik}
\end{figure}

Dalam implementasi teknis menggunakan LangGraph, konsep teoritis ReAct yang diajukan oleh \textcite{yaoReActSynergizingReasoning2023} diterjemahkan ke dalam struktur graf berstatus (\textit{stateful graph}). Pada desain aslinya, pola ReAct diimplementasikan secara linier melalui teknik \textit{prompt engineering}, di mana LLM menghasilkan teks terstruktur berisi penalaran (\textit{thought}) dan tindakan (\textit{action}) yang kemudian diurai oleh lingkungan eksternal. Di lain sisi, LangGraph memodelkan agen sebagai sistem yang terdiri dari node (simpul) dan edge (sisi). Struktur ini memetakan siklus Thought-Action-Observation ke dalam dua komponen utama, yaitu node Model yang berisi agen dan node Tools yang berisi alat. Node Model bertanggung jawab untuk menentukan apakah suatu tugas memerlukan penggunaan alat berdasarkan konteks yang diberikan. Jika alat diperlukan, LLM akan menghasilkan permintaan pemanggilan alat (\textit{tool call}) yang valid secara skema untuk dieksekusi secara otomatis oleh node Tools. Hasil eksekusi tersebut dikembalikan ke node Model sebagai observasi untuk dianalisis lebih lanjut hingga solusi akhir tercapai. Status (\textit{state}) agen dikelola sebagai objek yang menyimpan riwayat pesan, yang berfungsi sebagai memori jangka pendek atau \textit{scratchpad}. Alur kerja dimulai dari node Model, di mana LLM memproses input dan menentukan langkah selanjutnya. Jika model memutuskan untuk menggunakan alat, alur dialihkan melalui conditional edge menuju node Tools. Setelah alat dieksekusi, hasilnya dicatat ke dalam state sebagai observasi, dan kendali dikembalikan ke node Model untuk mengevaluasi informasi baru tersebut. Proses iteratif ini berlanjut hingga model menilai bahwa informasi yang dikumpulkan telah memadai untuk memberikan jawaban akhir (END). Visualisasi perbandingan antara desain teoritis ReAct dengan implementasi berbasis graf pada LangGraph disajikan pada Gambar \ref{fig:react_langgraph_comparison}.

\begin{figure}[H]
    \centering
    \includegraphics[width=1\textwidth]{bab3/reactlanggraph.png}
    \caption[Perbandingan antara (a) konsep teoritis ReAct dan (b) implementasi berbasis graf pada LangGraph]{Perbandingan antara (a) konsep teoritis ReAct oleh \textcite{yaoReActSynergizingReasoning2023} dan (b) implementasi berbasis graf pada LangGraph}
    \label{fig:react_langgraph_comparison}
\end{figure}

\section{Perancangan dan Pengembangan}

\subsection{Perancangan dan Pengembangan Agen LLM}
Bagian ini menjelaskan spesifikasi dan konfigurasi agen LLM yang dikembangkan untuk mengimplementasikan pola ReAct pada sistem identifikasi penyakit tanaman.

\subsubsection{Konfigurasi}
Agen dalam penelitian ini digerakkan oleh Gemini 3 Flash sebagai LLM utama. Pemilihan model ini didasarkan pada karakteristik model yang menawarkan kecepatan inferensi tinggi dan kemampuan penalaran (\textit{reasoning}) serta penggunaan alat (\textit{tool-use}) yang sangat baik untuk tugas-tugas agentik. Konfigurasi model diatur dengan nilai \textit{temperature} 0 guna memastikan luaran yang bersifat deterministik dan konsisten, sehingga meminimalkan variasi jawaban yang tidak relevan dalam konteks teknis identifikasi penyakit penyakit.
%TODO: pastiin kata diagnosis penyakit diganti ke identifasi penyakti
Melalui implementasi arsitekur LangGraph, agen memiliki memori jangka pendek (\textit{short-term memory}) yang dikelola dikelola melalui objek \textit{State} (status) yang bersifat \textit{stateful}. Status ini berfungsi sebagai buku catatan (\textit{scratchpad}) yang menyimpan seluruh riwayat interaksi dan informasi yang dikumpulkan selama satu sesi tugas berlangsung. Dengan menyimpan status secara terstruktur, agen memiliki konteks terhadap langkah-langkah yang telah diambil sebelumnya, misalnya hasil deteksi objek atau temuan dari pencarian internet, yang kemudian digunakan sebagai dasar untuk menentukan langkah tindakan selanjutnya hingga solusi akhir tercapai. Tabel \ref{tab:state_components} menyajikan rincian variabel status yang dikelola dalam sistem.

\begin{table}[H]
    \centering
    \caption{Komponen Status (State) pada Agen LLM}
    \label{tab:state_components}
    \small
    \setstretch{1.0}
    \begin{tabular}{|l|l|p{4cm}|}
    \hline
    Variabel Status & Tipe Data & Deskripsi \\ \hline
    \texttt{messages} & \texttt{list[AnyMessage]} & Menyimpan riwayat seluruh pesan interaksi (percakapan) antara pengguna dan agen sebagai memori jangka pendek (\textit{scratchpad}). \\ \hline
    \texttt{images} & \texttt{list[dict]} & Menyimpan daftar gambar yang diunggah beserta metadata terkait untuk mendukung penanganan beberapa citra sekaligus. \\ \hline
    \texttt{detections} & \texttt{list[dict]} & Menyimpan data hasil deteksi objek yang mencakup label, skor kepercayaan, dan koordinat kotak pembatas (\textit{bounding box}). \\ \hline
    \texttt{plant\_disease\_classifications} & \texttt{list[dict]} & Menyimpan hasil identifikasi penyakit tanaman yang mencakup label prediksi, tingkat kepercayaan, dan metadata pendukung lainnya. \\ \hline
    \texttt{visualization\_urls} & \texttt{list[str]} & Menyimpan daftar tautan publik dari citra visualisasi hasil deteksi yang telah diunggah ke penyimpanan Cloudflare R2. \\ \hline
    \texttt{current\_image\_url} & \texttt{str} & Menyimpan URL dari gambar yang sedang menjadi fokus analisis aktif dalam proses iterasi berjalan. \\ \hline
    \texttt{jump\_to} & \texttt{JumpTo} & Atribut kontrol internal untuk mengelola navigasi alur kerja (\textit{routing}) antar simpul (\textit{node}) dalam arsitektur graf LangGraph. \\ \hline
    \end{tabular}
\end{table}

Untuk memastikan stabilitas, keamanan, dan fungsionalitas yang optimal, agen dilengkapi dengan serangkaian \textit{middleware}. \textit{Middleware} dalam arsitektur ini berfungsi sebagai lapisan perantara yang membungkus pemanggilan model maupun alat (\textit{tools}) untuk pemrosesan tambahan secara otomatis sebelum dan sesudah eksekusi dilakukan. Tabel \ref{tab:middleware_components} merinci komponen \textit{middleware} yang diintegrasikan ke dalam sistem.

\begin{table}[H]
    \centering
    \caption{Komponen Middleware pada Agen LLM}
    \label{tab:middleware_components}
    \small
    \setstretch{1.0}
    \begin{tabular}{|l|p{7cm}|}
    \hline
    Middleware & Deskripsi Fungsi \\ \hline
    \texttt{get\_system\_prompt} & Menyediakan instruksi dasar, batasan peran, dan panduan penggunaan alat kepada agen. \\ \hline
    \texttt{handle\_tool\_errors} & Menangkap kegagalan eksekusi pada \textit{tools} dan mengembalikan pesan kesalahan yang informatif agar agen dapat melakukan koreksi secara mandiri. \\ \hline
    \texttt{model\_retry\_middleware} & Melakukan percobaan ulang otomatis jika terjadi kegagalan koneksi atau batasan kuota (\textit{rate limit}) pada API model. \\ \hline
    \texttt{tool\_retry\_middleware} & Menangani percobaan ulang jika eksekusi alat eksternal mengalami gangguan teknis sementara. \\ \hline
    \texttt{model\_fallback\_middleware} & Mekanisme cadangan yang secara otomatis mengalihkan tugas ke model yang lebih tinggi, yakni Gemini 3 Pro apabila model utama gagal merespons. \\ \hline
    \end{tabular}
\end{table}

\subsubsection{\textit{Prompt Engineering}}
Desain instruksi (\textit{prompt engineering}) pada agen LLM dirancang dengan mengadopsi kerangka kerja CO-STAR (Context, Objective, Style, Tone, Audience, Response). Kerangka kerja ini digunakan untuk memastikan bahwa model mendapatkan konteks yang lengkap dan terstruktur guna menghasilkan respons yang akurat, relevan, dan konsisten \parencite{govtechsingaporeCOSTARPromptEngineering2024}. Gambar \ref{fig:costar_framework} menyajikan visualisasi dari kerangka kerja CO-STAR. CO-STAR merupakan akronim dari enam komponen utama yang dijelaskan sebagai berikut:
\begin{enumerate}[nosep]
    \item \textit{Context} (Konteks): Memberikan informasi latar belakang mengenai peran agen sebagai asisten identifikasi penyakit tanaman, batasan domain penelitian, serta kapabilitas sistem.
    \item \textit{Objective} (Tujuan): Menetapkan tujuan utama agen untuk menghasilkan identifikasi yang berbasis bukti, jujur terhadap ketidakpastian, dan memberikan panduan yang berlandaskan data.
    \item \textit{Style} (Gaya): Menentukan cara penulisan respons agar mengalir secara alami, eksplisit dalam penyebutan sumber bukti, serta ringkas tanpa pengulangan informasi.
    \item \textit{Tone} (Nada): Mengatur nada bicara agen agar bersifat praktis, tenang, dan profesional tanpa memberikan peringatan yang berlebihan.
    \item \textit{Audience} (Audiens): Mengidentifikasi target pengguna sistem yang mencakup petani, penghobi tanaman, hingga tenaga profesional di bidang pertanian dengan tingkat pemahaman teknis yang bervariasi.
    \item \textit{Response} (Respons): Menentukan struktur luaran yang diharapkan, mulai dari observasi visual, integrasi konteks pengguna, hingga hipotesis diagnosis dan saran tindak lanjut.
\end{enumerate}
%TODO: perbagus lagi diagram ini
\begin{figure}[H]
    \centering
    \begin{tikzpicture}[
        node distance=1.5cm,
        main/.style={circle, draw, fill=blue!10, align=center, font=\small},
        comp/.style={rectangle, draw, fill=green!10, align=center, font=\footnotesize, minimum width=2cm}
    ]
    \node (costar) [main, minimum size=2cm] {CO-STAR\\Prompt};

    \node (c) [comp, above=of costar] {Context};
    \node (o) [comp, above right=1cm and 1.5cm of costar] {Objective};
    \node (s) [comp, below right=1cm and 1.5cm of costar] {Style};
    \node (t) [comp, below=of costar] {Tone};
    \node (a) [comp, below left=1cm and 1.5cm of costar] {Audience};
    \node (r) [comp, above left=1cm and 1.5cm of costar] {Response};

    \draw[-stealth] (c) -- (costar);
    \draw[-stealth] (o) -- (costar);
    \draw[-stealth] (s) -- (costar);
    \draw[-stealth] (t) -- (costar);
    \draw[-stealth] (a) -- (costar);
    \draw[-stealth] (r) -- (costar);
    \end{tikzpicture}
    \caption{Diagram Kerangka Kerja CO-STAR}
    \label{fig:costar_framework}
\end{figure}

Selain komponen standar CO-STAR, desain prompt pada sistem ini juga dilengkapi dengan bagian penjelasan alat (\textit{Tools Reference}). Bagian ini berfungsi sebagai dokumentasi teknis yang menjelaskan fungsionalitas, parameter, dan skenario penggunaan dari setiap alat yang terintegrasi, seperti pencarian web, pencarian basis pengetahuan, dan model deteksi visual. Dengan adanya penjelasan alat yang komprehensif, agen dapat memahami cara berinteraksi dengan lingkungan eksternal secara tepat.

\subsection{Perancangan dan Pengembangan Modul Deteksi Objek Tanaman}
Modul deteksi objek tanaman dirancang sebagai tahap awal dalam alur kerja identifikasi penyakit untuk melakukan lokalisasi terhadap bagian tanaman yang relevan, seperti daun, buah, atau batang. Modul ini memungkinkan agen untuk memperoleh kotak pembatas (\textit{bounding box}) dari objek yang terdeteksi, yang kemudian digunakan oleh modul pencarian gambar–teks untuk melakukan pencarian berbasis wilayah (\textit{region-based search}). Pendekatan ini bertujuan untuk mereduksi derau latar belakang (\textit{background noise}) serta memitigasi gangguan dari objek non-relevan lainnya dalam satu citra, sehingga meningkatkan akurasi pengambilan informasi.  Alur kerja dari modul deteksi objek tanaman divisualisasikan pada Gambar \ref{fig:plant_detection_workflow}.

Implementasi modul ini menggunakan pendekatan hibrida yang menggabungkan dua alat dengan karakteristik yang berbeda, yakni alat deteksi \textit{closed-set} dan alat deteksi \textit{open-vocabulary}. Alat deteksi \textit{closed-set} menggunakan arsitektur YOLOv11 dari \textcite{khanamYOLOv11OverviewKey2024} yang dioptimalkan khusus untuk deteksi daun, mengingat daun merupakan bagian tanaman yang paling dominan dalam memperlihatkan gejala penyakit. Sementara itu, alat deteksi \textit{open-vocabulary} menggunakan model OWLv2 dari \textcite{mindererScalingOpenVocabularyObject2024} untuk mendeteksi bagian tanaman lain melalui perintah teks (\textit{text prompt}) secara fleksibel. YOLOv11 menawarkan kecepatan inferensi dan akurasi yang sangat tinggi, namun untuk kelas terbatas yang spesifik berupa daun, sedangkan OWLv2 memberikan fleksibilitas deteksi objek kelas selain daun, seperti buah atau batang, namun dengan beban komputasi yang lebih besar.

\begin{figure}[H]
    \centering
    \includegraphics[width=1\textwidth]{bab3/moduldeteksitanaman.png}
    \caption{Alur Kerja Modul Deteksi Objek Tanaman}
    \label{fig:plant_detection_workflow}
\end{figure}

% TODO(Human/AI): Buat visualisasi kerja modul ini

\subsubsection{Alat Deteksi Closed-Set}
Alat deteksi \textit{closed-set} dalam modul ini diimplementasikan menggunakan arsitektur YOLOv11. Pemilihan YOLOv11 didasarkan pada performanya yang unggul dalam menyeimbangkan antara kecepatan inferensi dan akurasi deteksi dibandingkan versi sebelumnya. Model ini dilatih secara khusus untuk mengenali objek daun di lapangan.

\subsubsection{Persiapan Data}
Dataset yang digunakan untuk melatih model deteksi daun diadopsi dari dataset PlantDoc dari \textcite{singhPlantDocDatasetVisual2020}. Dataset ini dipilih karena mencakup variasi spesies tanaman dan kondisi penyakit yang luas dalam lingkungan alami. Dalam penelitian ini, dilakukan prapemrosesan berupa penyatuan seluruh anotasi kelas yang ada menjadi satu kelas tunggal, yakni "\textit{leaf}" (daun). Hal ini dilakukan untuk memfokuskan model pada tugas lokalisasi daun sebagai \textit{Region of Interest} (RoI) utama. Distribusi data awal dari dataset PlantDoc sebelum penyatuan kelas disajikan pada Gambar \ref{fig:plantdoc_distribution} untuk memberikan gambaran keragaman data yang digunakan.

\begin{figure}[H]
    \centering
    \includegraphics[width=1\textwidth]{bab3/plantdoc_distribution.png}
    \caption{Distribusi Kelas pada Dataset PlantDoc dari \textcite{singhPlantDocDatasetVisual2020}}
    \label{fig:plantdoc_distribution}
\end{figure}

\paragraph{Pemodelan}
Pemodelan dilakukan dengan mengevaluasi tiga variasi skala model YOLOv11, yaitu YOLOv11n (nano), YOLOv11s (small), dan YOLOv11m (medium). Ketiga variasi model ini dipilih karena kesesuaiannya dengan target lingkungan inferensi berbasis CPU serta kebutuhan sumber daya komputasi pelatihan yang masih dapat diakomodasi oleh perangkat keras yang tersedia dalam penelitian ini. Ketiga varian ini dibandingkan untuk menentukan model yang paling optimal untuk diintegrasikan ke dalam sistem agentik. Seluruh proses pelatihan dilakukan menggunakan Ultralytics dengan konfigurasi hiperparameter bawaan, kecuali yang dirinci pada Tabel \ref{tab:yolo_hyperparameters}.

\begin{table}[H]
\centering
    \caption{Konfigurasi Hiperparameter Pelatihan YOLOv11}
    \label{tab:yolo_hyperparameters}
    \centering
    \setstretch{1.0}
    \begin{tabular}{|l|l|p{6cm}|}
\hline
    Parameter & Nilai & Keterangan \\ \hline
    Epoch & 240 & Jumlah iterasi pelatihan penuh. \\ \hline
    Input Size & 416 x 416 & Resolusi citra masukan sesuai anotasi PlantDoc. \\ \hline
    Random Seed & 42 & Menjamin keterulangan hasil eksperimen. \\ \hline
    Batch Size & 16 & Jumlah sampel per sub-iterasi. \\ \hline
    Patience & 20 & Batas iterasi tanpa peningkatan untuk \textit{early stopping}. \\ \hline
    \end{tabular}
\end{table}

Setelah proses pelatihan selesai, model diekspor ke format ONNX (\textit{Open Neural Network Exchange}) untuk meningkatkan portabilitas dan efisiensi saat dijalankan di lingkungan produksi. Proses ekspor dilakukan dengan pengaturan resolusi sesuai ukuran masukan (\textit{input size}) pada perangkat CPU guna memastikan kompatibilitas dengan keterbatasan sumber daya pada lingkungan server.

\paragraph{Evaluasi}
Evaluasi model deteksi dilakukan untuk mengukur ketepatan lokalisasi dan klasifikasi objek. Metrik yang digunakan meliputi \textit{Precision}, \textit{Recall}, \textit{F1-score}, serta \textit{mean Average Precision} (mAP) pada ambang batas IoU 50\% (mAP50) dan rentang 50--95\% (mAP50-95). Definisi metrik tersebut diuraikan sebagai berikut:

\begin{enumerate}[nosep]
    \item IoU (\textit{Intersection over Union}): Mengukur tingkat tumpang tindih antara kotak prediksi ($BB_{pred}$) dan \textit{ground truth} ($BB_{gt}$). Secara praktis, metrik ini menunjukkan seberapa presisi model dalam melokalisasi posisi objek daun yang sebenarnya.
    \[IoU = \frac{Area(BB_{pred} \cap BB_{gt})}{Area(BB_{pred} \cup BB_{gt})}\]

    \item Precision: Rasio prediksi positif yang benar terhadap total prediksi positif. Metrik ini mengindikasikan tingkat kepercayaan terhadap prediksi model, yakni seberapa sering model benar ketika menyatakan mendeteksi daun dan meminimalkan deteksi palsu.
    \[Precision = \frac{TP}{TP + FP}\]

    \item Recall: Rasio prediksi positif yang benar terhadap total sampel positif aktual. Metrik ini mengukur sensitivitas model, yakni kemampuannya untuk menemukan seluruh objek daun yang ada dalam citra dan meminimalkan objek yang terlewat.
    \[Recall = \frac{TP}{TP + FN}\]

    \item F1-score: Rata-rata harmonis antara \textit{precision} dan \textit{recall}. Metrik ini memberikan gambaran keseimbangan performa model, memastikan model tidak hanya berhati-hati dalam mendeteksi tetapi juga mampu menemukan objek sebanyak mungkin.
    \[F1 = 2 \cdot \frac{Precision \cdot Recall}{Precision + Recall}\]

    \item mAP (\textit{mean Average Precision}): Rata-rata \textit{Average Precision} (AP) di seluruh kelas. Metrik ini merupakan standar evaluasi menyeluruh yang merepresentasikan keandalan model pada berbagai tingkat ambang batas deteksi.
    \[mAP = \frac{1}{N} \sum_{i=1}^{N} AP_i\]
\end{enumerate}

Selain metrik akurasi, dilakukan pula evaluasi terhadap performa waktu atau latensi inferensi menggunakan metrik persentil ke-50 (p50) dan persentil ke-95 (p95). Misalkan $T = \{t_1, t_2, \dots, t_n\}$ adalah himpunan waktu latensi yang telah diurutkan secara menaik, maka definisi dan kegunaan praktis dari kedua metrik tersebut adalah:

\begin{enumerate}[nosep]
    \item p50 (Median): Merepresentasikan nilai tengah dari distribusi latensi. Metrik ini menggambarkan kecepatan respons tipikal yang akan dialami oleh mayoritas pengguna dalam penggunaan normal.
    \[ p50 = \begin{cases} t_{\frac{n+1}{2}} & \text{jika } n \text{ ganjil} \\ \frac{1}{2} (t_{\frac{n}{2}} + t_{\frac{n}{2}+1}) & \text{jika } n \text{ genap} \end{cases} \]

    \item p95 (Persentil ke-95): Merepresentasikan ambang batas waktu di mana 95\% dari total sampel inferensi lebih cepat dari nilai tersebut. Metrik ini  z untuk mengukur performa dalam skenario terburuk (\textit{tail latency}), memastikan sistem tetap responsif bahkan saat memproses citra yang kompleks atau dalam kondisi beban tinggi.
    \[ p95 = t_{\lceil 0.95 \cdot n \rceil} \]
\end{enumerate}

Setelah tahap evaluasi, dilakukan \textit{benchmarking} dengan menganalisis hubungan antara mAP dan latensi melalui grafik \textit{Pareto frontier}. Grafik \textit{Pareto frontier} digunakan untuk mengidentifikasi sekumpulan model yang memberikan kombinasi performa terbaik tanpa ada model lain yang lebih unggul di kedua aspek secara bersamaan \parencite{shakyaSmallObjectDetectionEdge2026}. Visualisasi contoh grafik \textit{Pareto frontier} disajikan pada Gambar \ref{fig:pareto_curve}. Analisis ini bertujuan untuk mengidentifikasi varian model yang menawarkan latensi paling minimal dengan tetap mempertahankan performa deteksi yang optimal.

\begin{figure}[H]
    \centering
    \includegraphics[width=0.9\textwidth]{bab3/contohparetofrontier.png}
    \caption{Contoh Grafik \textit{Pareto Frontier} mAP vs Latensi \parencite{shakyaSmallObjectDetectionEdge2026}}
    \label{fig:pareto_curve}
\end{figure}

\subsubsection{Alat Deteksi Open-Vocabulary}
Alat deteksi \textit{open-vocabulary} menggunakan model OWLv2 yang dikembangkan oleh Google \parencite{mindererScalingOpenVocabularyObject2024}. Sistem menggunakan varian model \texttt{google/owlv2-base-patch16} yang telah dikonversi ke format ONNX. Model ini mengintegrasikan \textit{vision transformer} (ViT) dengan \textit{text encoder} untuk melakukan deteksi objek berdasarkan deskripsi tekstual secara \textit{zero-shot}. Gambar \ref{fig:arsitektur_owlv2} menyajikan visualisasi arsitektur dari model OWLv2.

\begin{figure}[H]
    \centering
    \includegraphics[width=1\textwidth]{bab3/arsitekturalatowlv2.png}
    \caption{Arsitektur Model OWLv2 yang Digunakan}
    \label{fig:arsitektur_owlv2}
\end{figure}

\subsection{Perancangan dan Pengembangan Modul Pencarian Gambar-Teks Penyakit}
Modul pencarian gambar-teks penyakit dirancang sebagai komponen utama sistem untuk mengidentifikasi gejala penyakit tanaman melalui pendekatan pengambilan informasi multimodal. Implementasi modul ini menggunakan teknologi Qdrant yang mendukung penyimpanan dan pencarian vektor bernama jamak (\textit{multi-named vectors}) secara efisien. Melalui fitur ini, dataset yang berisi pasangan gambar dan teks dapat disematkan (\textit{embedded}) ke dalam satu koleksi Qdrant, di mana satu poin data menampung dua vektor sekaligus, yaitu vektor gambar dan vektor teks. Representasi fitur diekstraksi menggunakan model \textit{foundation} SCOLD yang dikembangkan oleh \textcite{quocVisionLanguageFoundationModel2025} dengan arsitektur RoBERTa dari \textcite{liuRoBERTaRobustlyOptimized2019} sebagai \textit{encoder} teks dan arsitektur Swin-T dari \textcite{liuSwinTransformerHierarchical2021} sebagai \textit{encoder} gambar. Rincian lebih lanjut mengenai skema ini akan dibahas pada subbab berikutnya. Alur kerja dari modul pencarian gambar-teks divisualisasikan pada Gambar \ref{fig:image_text_search_workflow}.

\begin{figure}[H]
    \centering
    \includegraphics[width=0.9\textwidth]{bab3/modulgambarteks.png}
    \caption{Alur Kerja Modul Pencarian Gambar-Teks Penyakit Tanaman}
    \label{fig:image_text_search_workflow}
\end{figure}


\subsubsection{Persiapan Data}
Dataset PlantWild yang disusun oleh \textcite{weiBenchmarkingInthewildMultimodal2024} digunakan sebagai basis indeks pencarian pada modul ini. PlantWild mencakup 184.800 citra yang terdistribusi ke dalam 89 kelas kondisi kesehatan dan penyakit pada 13 spesies tanaman utama. Di mana, terdapat 56 kelas penyakit dan 33 kelas kondisi sehat. Gambar \ref{fig:plantwild_distribution} menyajikan distribusi kelas pada dataset PlantWild.

\begin{figure}[H]
    \centering
    \includegraphics[width=1\textwidth]{bab3/distribusiplantwild.png}
    \caption{Distribusi Kelas pada Dataset PlantWild dari \textcite{weiBenchmarkingInthewildMultimodal2024}}
    \label{fig:plantwild_distribution}
\end{figure}

\subsubsection{\textit{Indexing}}
Dalam Qdrant, data dikelola dalam unit entitas yang disebut sebagai Poin (\textit{Points}). Setiap Poin terdiri dari tiga komponen utama, yaitu ID unik, vektor yang dapat berupa vektor tunggal atau jamak, dan \textit{payload} opsional untuk menyimpan metadata. Pada sistem ini, setiap entitas penyakit tanaman direpresentasikan sebagai Poin yang memiliki dua vektor padat (\textit{dense vectors}) bernama, yaitu vektor teks dan vektor citra. Citra referensi asli disimpan secara terpisah pada layanan penyimpanan objek Cloudflare R2 untuk efisiensi, sementara URL publiknya disimpan dalam \textit{payload}.

Tabel \ref{tab:qdrant_schema} menyajikan rincian skema \textit{Points} dan struktur \textit{payload} yang diimplementasikan. Skema ini dirancang untuk mendukung pencarian multimodal sekaligus penyaringan metadata yang presisi.

%TODO: perlu diperbagus tabelnya
\begin{table}[H]
    \centering
    \caption{Skema Poin (Vektor dan Payload) pada Indeks Qdrant}
    \label{tab:qdrant_schema}
    \small
    \setstretch{1.0}
    \resizebox{\textwidth}{!}{%
    \begin{tabular}{|l|l|l|p{7cm}|}
    \hline
    \multicolumn{1}{|c|}{\textbf{Komponen}} & \multicolumn{1}{c|}{\textbf{Kunci}} & \multicolumn{1}{c|}{\textbf{Tipe}} & \multicolumn{1}{c|}{\textbf{Deskripsi}} \\ \hline
    \textbf{Vektor} & \texttt{text} & 512 & Representasi semantik deskripsi gejala (RoBERTa). \\ \hline
    & \texttt{image} & 512 & Representasi fitur visual citra tanaman (Swin-T/SCOLD). \\ \hline
    \textbf{Payload} & \texttt{plant\_name} & String & Nama spesies tanaman. Diindeks secara \textit{full-text} untuk filtrasi. \\ \hline
    & \texttt{label} & String & Nama spesifik penyakit. Diindeks secara \textit{full-text}. \\ \hline
    & \texttt{caption} & String & Deskripsi tekstual lengkap mengenai gejala. \\ \hline
    & \texttt{image\_url} & String & Tautan publik citra di Cloudflare R2. \\ \hline
    & \texttt{source\_id} & Integer & ID referensi internal. \\ \hline
    & \texttt{filename} & String & Nama fail asli citra. \\ \hline
    \end{tabular}%
    }
\end{table}

Fitur penting dalam implementasi ini adalah penerapan indeks teks penuh (\textit{full-text index}) pada kolom \textit{payload} \texttt{plant\_name} dan \texttt{label}. Berbeda dengan pencocokan kata kunci eksak (\textit{keyword match}), indeks teks penuh memungkinkan agen untuk melakukan penyaringan (\textit{filtering}) yang lebih fleksibel, seperti pencocokan berdasarkan token kata (\textit{token-based matching}) dan ketidakpekaan terhadap huruf besar-kecil (\textit{case-insensitive}). Fitur ini secara signifikan meningkatkan efisiensi pencarian agen dengan membatasi ruang pencarian vektor hanya pada subset data yang relevan dengan konteks tanaman yang sedang dianalisis. Sebagai ilustrasi, ketika agen mendeteksi bahwa tanaman yang dianalisis adalah tomat, agen dapat menerapkan filter pada \texttt{plant\_name} untuk hanya mencari di antara data penyakit tomat, sehingga hasil yang diperoleh lebih akurat dan relevan. Visualisasi mekanisme penyaringan menggunakan indeks teks penuh pada \textit{payload} disajikan pada Gambar \ref{fig:payload_filtering}.

%TODO: ini gambar masih placeholder, sesuain konteksnya ke tanaman!
\begin{figure}[H]
    \centering
    \includegraphics[width=1\textwidth]{bab3/placeholderfulltextindexpayloadfield.png}
    \caption{Visualisasi mekanisme penyaringan menggunakan \textit{full-text index} pada \textit{payload} Qdrant}
    \label{fig:payload_filtering}
\end{figure}

\subsubsection{Strategi Pencarian}
Implementasi modul ini mendukung empat strategi pencarian multimodal yang memanfaatkan kemampuan representasi vektor ganda (teks dan citra) pada indeks Qdrant. Pemilihan strategi dilakukan secara dinamis oleh agen berdasarkan ketersediaan input dari pengguna (teks gejala atau citra tanaman) serta tujuan pencarian. Fleksibilitas ini dimungkinkan oleh arsitektur model SCOLD yang menyelaraskan (\textit{align}) representasi fitur teks dan citra ke dalam ruang laten yang sama, memungkinkan operasi pencarian lintas modal (\textit{cross-modal}). Keempat strategi tersebut diilustrasikan pada Gambar \ref{fig:multimodal_search_strategies} dan dijelaskan sebagai berikut:

%TODO: INI GAMBAR MASIH PLACEHOLDER, BUAT YANG LEBIH BAGUS DAN SESUAI KONTEX TANAMAN
\begin{figure}[H]
    \centering
    \begin{tikzpicture}[
        node distance=0.5cm and 0.8cm,
        font=\scriptsize,
        box/.style={rectangle, draw, rounded corners, align=center, minimum width=1.8cm, minimum height=0.8cm, fill=white},
        db/.style={cylinder, shape border rotate=90, draw, aspect=0.25, minimum width=1.2cm, minimum height=1.2cm, align=center, fill=gray!10},
        arrow/.style={-stealth, thick},
        label/.style={font=\bfseries, align=center}
    ]

    % --- (a) Text-to-Text ---
    \node (t2t_input) [box, fill=blue!10] {Input: Teks\\(Gejala)};
    \node (t2t_encoder) [box, right=of t2t_input] {Text Encoder\\(RoBERTa)};
    \node (t2t_vector) [box, right=of t2t_encoder, fill=yellow!10] {Vektor\\Teks};
    \node (t2t_db) [db, right=of t2t_vector] {Qdrant\\\tiny(Vec: Text)};

    \draw [arrow] (t2t_input) -- (t2t_encoder);
    \draw [arrow] (t2t_encoder) -- (t2t_vector);
    \draw [arrow] (t2t_vector) -- (t2t_db);

    \node [below=0.1cm of t2t_encoder] {(a) Text-to-Text};

    % --- (b) Text-to-Image ---
    \node (t2i_input) [box, fill=blue!10, below=1.5cm of t2t_input] {Input: Teks\\(Gejala)};
    \node (t2i_encoder) [box, right=of t2i_input] {Text Encoder\\(RoBERTa)};
    \node (t2i_vector) [box, right=of t2i_encoder, fill=yellow!10] {Vektor\\Teks};
    \node (t2i_db) [db, right=of t2i_vector] {Qdrant\\\tiny(Vec: Image)};

    \draw [arrow] (t2i_input) -- (t2i_encoder);
    \draw [arrow] (t2i_encoder) -- (t2i_vector);
    \draw [arrow] (t2i_vector) -- (t2i_db);

    \node [below=0.1cm of t2i_encoder] {(b) Text-to-Image};

    % --- (c) Image-to-Image ---
    \node (i2i_input) [box, fill=green!10, below=1.5cm of t2i_input] {Input: Citra\\(Daun)};
    \node (i2i_encoder) [box, right=of i2i_input] {Image Encoder\\(SCOLD)};
    \node (i2i_vector) [box, right=of i2i_encoder, fill=orange!10] {Vektor\\Citra};
    \node (i2i_db) [db, right=of i2i_vector] {Qdrant\\\tiny(Vec: Image)};

    \draw [arrow] (i2i_input) -- (i2i_encoder);
    \draw [arrow] (i2i_encoder) -- (i2i_vector);
    \draw [arrow] (i2i_vector) -- (i2i_db);

    \node [below=0.1cm of i2i_encoder] {(c) Image-to-Image};

    % --- (d) Image-to-Text ---
    \node (i2t_input) [box, fill=green!10, below=1.5cm of i2i_input] {Input: Citra\\(Daun)};
    \node (i2t_encoder) [box, right=of i2t_input] {Image Encoder\\(SCOLD)};
    \node (i2t_vector) [box, right=of i2t_encoder, fill=orange!10] {Vektor\\Citra};
    \node (i2t_db) [db, right=of i2t_vector] {Qdrant\\\tiny(Vec: Text)};

    \draw [arrow] (i2t_input) -- (i2t_encoder);
    \draw [arrow] (i2t_encoder) -- (i2t_vector);
    \draw [arrow] (i2t_vector) -- (i2t_db);

    \node [below=0.1cm of i2t_encoder] {(d) Image-to-Text};

    \end{tikzpicture}
    \caption{Visualisasi Empat Strategi Pencarian Multimodal: (a) Text-to-Text, (b) Text-to-Image, (c) Image-to-Image, dan (d) Image-to-Text}
    \label{fig:multimodal_search_strategies}
\end{figure}

Penjelasan rinci mengenai mekanisme kerja dari masing-masing strategi pencarian tersebut diuraikan sebagai berikut:

\begin{enumerate}[nosep]
    \item Text-to-Text (Pencarian Semantik Teks) \\
    Mode ini digunakan ketika pengguna memberikan deskripsi gejala berupa teks. Input diproses oleh \textit{tokenizer} dan \textit{text encoder} RoBERTa untuk menghasilkan vektor teks. Vektor ini kemudian dicocokkan dengan koleksi vektor teks (\texttt{text}) dalam indeks Qdrant yang merepresentasikan \textit{caption} penyakit. Strategi ini efektif untuk menemukan diagnosis berdasarkan deskripsi linguistik gejala.

    \item Text-to-Image (Pencarian Lintas Modal Teks-ke-Citra) \\
    Mode ini memungkinkan pencarian citra referensi penyakit hanya dengan menggunakan deskripsi teks. Vektor teks hasil enkripsi digunakan untuk mencari tetangga terdekat pada ruang vektor citra (\texttt{image}). Kemampuan ini dimungkinkan karena model SCOLD telah dilatih untuk meminimalkan jarak antara pasangan teks dan citra yang relevan dalam ruang vektor (\textit{embedding space}).

    \item Image-to-Image (Pencarian Kemiripan Visual) \\
    Mode ini adalah metode standar untuk identifikasi berbasis citra. Citra masukan dari pengguna (baik citra utuh maupun hasil potongan/crop deteksi objek) diproses oleh \textit{image encoder} (Swin-Transformer basis SCOLD) menjadi vektor citra. Vektor ini kemudian dibandingkan dengan jutaan vektor citra (\texttt{image}) dalam basis data untuk menemukan penyakit dengan karakteristik visual paling serupa.

    \item Image-to-Text (Pencarian Lintas Modal Citra-ke-Teks) \\
    Mode ini membalikkan pendekatan lintas modal, di mana vektor citra masukan digunakan untuk mencari kecocokan pada ruang vektor teks (\texttt{text}). Strategi ini berguna untuk mendapatkan deskripsi tekstual atau label penyakit yang paling relevan secara langsung dari fitur visual, tanpa harus bergantung pada metadata citra referensi.
\end{enumerate}


\subsection{Perancangan dan Pengembangan Modul Pencarian Informasi}
Modul pencarian informasi dirancang untuk menyediakan konteks eksternal yang faktual bagi agen dalam menjawab pertanyaan pengguna atau memverifikasi hipotesis diagnosis. Agen diinstruksikan untuk memanfaatkan modul ini ketika pengguna meminta panduan manajemen penyakit atau perawatan preventif tanaman. Selain itu, modul ini berfungsi sebagai mekanisme cadangan (\textit{fallback}) faktual terakhir untuk membantu identifikasi apabila modul pencarian gambar-teks penyakit tanaman gagal memberikan hasil yang relevan, misalnya karena penyakit yang diidentifikasi tidak tersedia dalam indeks basis data vektor. Modul ini mengimplementasikan mekanisme \textit{Retrieval-Augmented Generation} (RAG) yang terdiri dari dua alat pencarian, yaitu alat pencarian pada basis data pengetahuan internal dan alat pencarian internet menggunakan Tavily API.

Basis data pengetahuan internal dibangun dari korpus dokumen terkurasi yang berfokus pada penyakit tanaman, yang disimpan dalam basis data vektor Qdrant untuk efisiensi pencarian. Dalam perancangannya, kinerja pencarian pada basis data pengetahuan dioptimalkan melalui pengujian strategi \textit{hybrid search} dan penerapan mekanisme \textit{reranking} untuk meningkatkan efesiensi pencarian. Meskipun basis data pengetahuan menawarkan informasi yang lebih terpercaya, cakupannya terbatas pada spesies tanaman yang telah diindeks. Oleh karena itu, pencarian internet melalui Tavily API diintegrasikan sebagai mekanisme cadangan (\textit{fallback}) untuk menangani kueri yang tidak dapat dijawab oleh data internal. Agen dirancang untuk memprioritaskan penggunaan pencarian basis data pengetahuan terlebih dahulu sebelum mengakses informasi dari internet. Alur kerja interaksi antara agen dan modul pencarian informasi divisualisasikan pada Gambar \ref{fig:info_search_workflow}.

\begin{figure}[H]
    \centering
    \begin{tikzpicture}[
        node distance=1.5cm,
        font=\small,
        start/.style={circle, draw, fill=blue!10, minimum size=1.5cm, align=center},
        process/.style={rectangle, draw, rounded corners, fill=white, minimum width=3cm, minimum height=1cm, align=center},
        decision/.style={diamond, draw, fill=yellow!10, aspect=2, inner sep=0pt, minimum width=2.5cm, align=center},
        arrow/.style={-stealth, thick},
        label/.style={font=\footnotesize, fill=white, inner sep=2pt}
    ]

    \node (agent) [start] {Agen\\ReAct};
    \node (kb) [process, right=1.5cm of agent] {Pencarian Basis\\Pengetahuan (Qdrant)};
    \node (check) [decision, right=1cm of kb] {Informasi\\Ditemukan?};
    \node (web) [process, below=of kb] {Pencarian Internet\\(Tavily API)};
    \node (end) [start, right=1.5cm of check, fill=green!10] {Konteks\\Informasi};

    \draw [arrow] (agent) -- node[above] {Prioritas 1} (kb);
    \draw [arrow] (kb) -- (check);
    \draw [arrow] (check) -- node[above] {Ya} (end);
    \draw [arrow] (check) |- node[near start, right] {Tidak} (web);
    \draw [arrow] (web) -| node[near start, above] {Fallback (Prioritas 2)} (end);

    \end{tikzpicture}
    \caption{Alur kerja modul pencarian informasi dengan prioritas pada basis data pengetahuan}
    \label{fig:info_search_workflow}
\end{figure}

\subsubsection{Alat Pencarian Basis Data Pengetahuan}
Implementasi alat pencarian basis data pengetahuan diawali dengan pemilihan komponen model yang optimal melalui proses \textit{benchmarking}. Berdasarkan studi literatur, terdapat perbedaan performa antara berbagai model \textit{embedding} yang tersedia saat ini. Mengacu pada \textit{Massive Text Embedding Benchmark} (MTEB) oleh \textcite{muennighoffMTEBMassiveText2023}, model \texttt{gemini-embedding-001} dari Google DeepMind tercatat sebagai model proprietari dengan peringkat teratas, mengungguli model \texttt{text-embedding-3-large} dari OpenAI yang berada di peringkat kelima. Peringkat ini pada MTEB ditampilkan pada Gambar \ref{fig:mteb_benchmark}.

\begin{figure}[H]
    \centering
    \includegraphics[width=0.8\textwidth]{bab3/mtebbenchmark.png}
    \caption{Peringkat \texttt{gemini-embedding-001} dan \texttt{text-embedding-3-large} pada MTEB oleh \textcite{muennighoffMTEBMassiveText2023}}
    \label{fig:mteb_benchmark}
\end{figure}

Akan tetapi, hasil yang berbeda ditunjukkan pada \textit{Semantic Alignment \& Generalization Evaluation} (SAGE) oleh \textcite{goelSAGERealisticBenchmark2025}. Pada tolok ukur ini, model \texttt{text-embedding-3-large} menempati posisi terbaik dengan skor keseluruhan SAGE sebesar 0,524, sedangkan \texttt{gemini-embedding-001} berada di peringkat ketiga dengan skor 0,504.

Mengingat adanya perbedaan hasil evaluasi tersebut, dilakukan \textit{benchmarking} internal menggunakan korpus data yang akan digunakan pada sistem untuk menentukan konfigurasi model yang paling sesuai dengan karakteristik data penyakit tanaman yang digunakan. Pengujian difokuskan pada konfigurasi \textit{hybrid search} yang menggabungkan pencarian vektor padat (\textit{dense vector}) dan vektor jarang (\textit{sparse vector}). Untuk komponen vektor jarang, digunakan dua pendekatan populer yaitu BM25 dari \textcite{robertsonProbabilisticRelevanceFramework2009} melalui Qdrant dan SPLADE dari \textcite{formalSPLADESparseLexical2021} melalui FastEmbed. Kombinasi pengujian mencakup empat variasi konfigurasi \textit{hybrid search}, yaitu (1) \texttt{text-embedding-3-large} dengan BM25, (2) \texttt{text-embedding-3-large} dengan SPLADE, (3) \texttt{gemini-embedding-001} dengan BM25, dan (4) \texttt{gemini-embedding-001} dengan SPLADE. Konfigurasi \textit{hybrid search} terbaik selanjutnya akan diuji lebih lanjut untuk mengevaluasi efektivitas komponen \textit{reranker} menggunakan model \texttt{jina-reranker-v3} dari Jinai AI. Model ini dipilih karena merupakan model \textit{cross-encoder} yang mendapatkan peringkat tertinggi pada MTEB. Hal ini diperlihatkan pada Gambar \ref{fig:mteb_reranker_benchmark}.

\begin{figure}[H]
    \centering
    \includegraphics[width=1\textwidth]{bab3/mteb_reranker_benchmark.png}
    \caption{Peringkat Model Reranker pada MTEB oleh \textcite{muennighoffMTEBMassiveText2023}}
    \label{fig:mteb_reranker_benchmark}
\end{figure}

Untuk pembanding, model reranker \texttt{rerank-2.5} dari Voyage AI dibandingkan dengan \texttt{jina-reranker-v3} dari Jina AI guna mendapatkan konfigurasi \textit{Retrieval-Augmented Generation} (RAG) yang paling optimal sebelum diimplementasikan ke sistem agentik. Alur pengujian komponen pencarian ini divisualisasikan pada Gambar \ref{fig:alur_pengujian_retrieval}.

\begin{figure}[H]
    \centering
    \begin{tikzpicture}[
        node distance=1.2cm,
        start/.style={rectangle, draw, rounded corners, fill=blue!10, align=center, minimum height=1cm, font=\small},
        process/.style={rectangle, draw, fill=white, align=center, minimum height=1cm, font=\small},
        arrow/.style={-stealth, thick}
    ]

    \node (step1) [start] {Uji 4 Konfigurasi Hybrid Search\\(OpenAI/Gemini + BM25/SPLADE)};
    \node (step2) [process, below=of step1] {Identifikasi Konfigurasi\\Hybrid Search Terbaik};
    \node (step3) [process, below=of step2] {Uji Efektivitas Reranking\\(Voyage AI vs Jina AI)};
    \node (step4) [start, below=of step3, fill=green!10] {Konfigurasi RAG\\Paling Optimal};

    \draw [arrow] (step1) -- (step2);
    \draw [arrow] (step2) -- (step3);
    \draw [arrow] (step3) -- (step4);

    \end{tikzpicture}
    \caption{Diagram Alur Pengujian Komponen Alat Pencarian Basis Data Pengetahuan}
    \label{fig:alur_pengujian_retrieval}
\end{figure}

\paragraph{Persiapan Data}
Persiapan data untuk basis pengetahuan dilakukan melalui serangkaian tahapan mulai dari pengumpulan data mentah melalui \textit{web scraping}, penyusunan korpus, hingga pembangkitan dataset evaluasi sintetis. Alur kerja lengkap persiapan data ini divisualisasikan pada Gambar \ref{fig:data_preparation_workflow}.

\begin{figure}[H]
    \centering
    \begin{tikzpicture}[
        node distance=1.0cm,
        font=\scriptsize,
        process/.style={rectangle, draw, rounded corners, fill=blue!5, align=center, minimum width=3cm, minimum height=0.8cm},
        data/.style={cylinder, shape border rotate=90, draw, aspect=0.25, fill=yellow!10, align=center, minimum width=2cm, minimum height=1cm},
        model/.style={ellipse, draw, dashed, fill=green!5, align=center, minimum width=2cm},
        arrow/.style={-stealth, thick}
    ]

    % Nodes
    \node (sources) [process, fill=white] {Sumber Data\\(PlantVillage \& Gardenology)};
    \node (crawl) [process, below=of sources] {Crawling URL\\(Playwright/Crawl4AI)};
    \node (scrape) [process, below=of crawl] {Scraping \& Ekstraksi\\(ScrapeGraphAI)};
    \node (llm_scrape) [model, right=0.5cm of scrape] {LLM:\\Gemini 2.5 Flash};

    \node (json) [data, below=of scrape] {Data JSON\\(222 Entri)};

    \node (chunk) [process, below=of json] {Semantic Chunking\\(LlamaIndex)};
    \node (embed) [model, right=0.5cm of chunk] {Embedding:\\Gemini Emb-001};

    \node (corpus) [data, below=of chunk] {Korpus Dokumen\\(2.469 Dokumen)};

    \node (gen_qa) [process, below=of corpus] {Pembangkitan QA\\(AutoRAG Strategy)};
    \node (llm_qa) [model, right=0.5cm of gen_qa] {LLM:\\Gemini 2.5 Flash};

    \node (dataset) [data, below=of gen_qa, fill=green!20] {Dataset QA\\(600 Sampel)};

    % Arrows
    \draw [arrow] (sources) -- (crawl);
    \draw [arrow] (crawl) -- (scrape);
    \draw [arrow] (scrape) -- (json);
    \draw [arrow] (json) -- (chunk);
    \draw [arrow] (chunk) -- (corpus);
    \draw [arrow] (corpus) -- (gen_qa);
    \draw [arrow] (gen_qa) -- (dataset);

    \draw [arrow, dashed] (llm_scrape) -- (scrape);
    \draw [arrow, dashed] (embed) -- (chunk);
    \draw [arrow, dashed] (llm_qa) -- (gen_qa);

    \end{tikzpicture}
    \caption{Alur Kerja Persiapan Data Dari Crawling hingga Pembangkitan Dataset QA}
    \label{fig:data_preparation_workflow}
\end{figure}

%TODO: ini mungkin perlu sitasi ajh webnya
Data mentah dikumpulkan dari dua sumber utama, yaitu PlantVillage (plantvillage.psu.edu) dan Gardenology (gardenology.org). PlantVillage dipilih karena reputasinya sebagai sumber informasi penyakit tanaman yang kredibel yang dikelola oleh Penn State University, sedangkan Gardenology dipilih untuk melengkapi informasi mengenai aspek budidaya dan propagasi tanaman yang kurang mendalam pada PlantVillage.

Proses pengumpulan data dilakukan menggunakan teknik \textit{web scraping} hibrida yang menggabungkan automasi peramban (\textit{browser automation}) dan ekstraksi berbasis LLM. Untuk penemuan tautan (\textit{crawling}), digunakan pustaka Playwright dari Python. Selanjutnya, ekstraksi konten dilakukan menggunakan pustaka ScrapeGraphAI. Pustaka ini memanfaatkan LLM untuk mengurai struktur HTML yang tidak beraturan menjadi format JSON yang terstruktur sesuai skema yang didefinisikan. Model Gemini 2.5 Flash dengan temperatur 0 dipilih sebagai otak ekstraksi karena memiliki jendela konteks (\textit{context window}) yang besar dan latensi rendah. Tabel \ref{tab:scraping_schema} menyajikan skema data yang digunakan untuk ekstraksi. Dari proses ini, berhasil dikumpulkan 222 entri data yang mencakup sekitar 200 spesies tanaman unik.

\begin{table}[H]
    \centering
    \caption{Skema Ekstraksi Data dari PlantVillage dan Gardenology}
    \label{tab:scraping_schema}
    \small
    \begin{tabular}{|l|l|p{7cm}|}
    \hline
    \textbf{Atribut} & \textbf{Tipe Data} & \textbf{Deskripsi} \\ \hline
    \texttt{plant} & String & Nama umum tanaman. \\ \hline
    \texttt{description} & String & Deskripsi umum karakteristik tanaman. \\ \hline
    \texttt{propagation} & String & Informasi budidaya tanaman. \\ \hline
    \texttt{pests\_and\_diseases} & List/String & Daftar penyakit beserta gejala, penyebab, dan manajemennya. \\ \hline
    \end{tabular}
\end{table}

Data JSON hasil \textit{scraping} kemudian diproses menjadi korpus dokumen menggunakan kerangka kerja LlamaIndex. Informasi dipisahkan menjadi dua kategori dokumen utama, yaitu (1) dokumen budidaya/propagasi dan (2) dokumen penyakit. Untuk menangani dokumen teks yang panjang, diterapkan teknik pemotongan semantik (\textit{semantic chunking}) menggunakan \texttt{SemanticSplitterNodeParser} dari LlamaIndex dengan model \textit{embedding} \texttt{gemini-embedding-001}. Pemotongan semantik bekerja dengan menghitung jarak kosinus (\textit{cosine distance}) antar representasi vektor dari kalimat-kalimat yang berurutan. Dalam implementasinya, ditetapkan parameter \textit{breakpoint percentile threshold} sebesar 95. Nilai ini menginstruksikan algoritma untuk menentukan ambang batas pemotongan berdasarkan persentil ke-95 dari distribusi jarak semantik dalam dokumen. Dengan kata lain, pemisahan dokumen (\textit{splitting}) hanya dilakukan pada titik-titik yang memiliki tingkat perbedaan semantik tertinggi (top 5\%), yang mengindikasikan adanya pergantian topik atau konteks yang signifikan. Pendekatan ini memastikan bahwa setiap potongan dokumen (\textit{chunk}) memiliki koherensi tematik yang utuh dan tidak terpotong di tengah penjelasan yang relevan. Proses ini menghasilkan total 2.469 dokumen yang siap diindeks.

Untuk mengevaluasi performa pada \textit{benchmarking}, dibuat dataset \textit{Question-Answering} (QA) sintetis sebanyak 600 sampel dari korpus yang telah disusun. Pembangkitan dilakukan menggunakan Gemini 2.5 Flash dengan temperatur 0 dan dalam bahasa Inggris. Pembangkitan dataset ini menggunakan dan mengadopsi metodologi yang disediakan AutoRAG dari \textcite{kimAutoRAGAutomatedFramework2024}. Dataset mencakup tiga jenis pertanyaan dengan tingkat kompleksitas berbeda, sebagaimana dirangkum dalam Tabel \ref{tab:qa_types}.

\begin{table}[H]
    \centering
    \caption{Jenis Pertanyaan pada Dataset Evaluasi QA Sintetis}
    \label{tab:qa_types}
    \small
    \begin{tabular}{|l|p{4cm}|p{6cm}|c|}
    \hline
    \textbf{Tipe Pertanyaan} & \textbf{Definisi} & \textbf{Contoh} & \textbf{Proporsi} \\ \hline
    \textit{Factoid Single Hop} & Pertanyaan yang mencari fakta singkat dan spesifik yang dapat diverifikasi. & "Apa patogen penyebab penyakit hawar daun pada tomat?" & 35\% \\ \hline
    \textit{Concept Completion} & Pertanyaan yang menanyakan esensi atau definisi dari suatu konsep. & "Jelaskan karakteristik dari gejala klorosis." & 35\% \\ \hline
    \textit{Two-Hop Incremental} & Pertanyaan multi-lompatan yang membutuhkan inferensi dari dua segmen informasi berbeda \parencite{hwangExplainableMultihopQuestion2024}. & "Bagaimana cara menangani penyakit yang disebabkan oleh jamur Fusarium pada tanaman pisang?" & 30\% \\ \hline
    \end{tabular}
\end{table}

\paragraph{Indexing}
Proses pengindeksan dilakukan dengan memanfaatkan fitur vektor bernama jamak (\textit{multi-named vectors}) pada Qdrant. Fitur ini memungkinkan penyimpanan beberapa representasi vektor sekaligus dalam satu entitas Poin (\textit{Point}) tunggal. Pendekatan ini diterapkan untuk memfasilitasi evaluasi strategi pencarian yang komprehensif dengan menyatukan seluruh representasi vektor dari berbagai model dalam satu koleksi terpadu. Setiap dokumen dalam korpus diubah menjadi Poin yang memiliki empat komponen vektor, yaitu dua vektor padat (\textit{dense vectors}) dan dua vektor jarang (\textit{sparse vectors}). Vektor padat dihasilkan oleh model \texttt{gemini-embedding-001} dan \texttt{text-embedding-3-large}, masing-masing dengan dimensi 3.072. Sementara itu, vektor jarang dihasilkan menggunakan model SPLADE++ dan algoritma BM25 yang terintegrasi secara internal di Qdrant.

Selain representasi vektor, setiap Poin juga menyimpan \textit{payload} yang berisi konten teks dokumen dan metadata terkait. Untuk mendukung efisiensi pengambilan informasi oleh agen, diterapkan indeks teks penuh (\textit{full-text index}) pada atribut metadata nama tanaman (\texttt{metadata.plant}) dan nama penyakit (\texttt{metadata.disease}). Penggunaan indeks teks penuh memungkinkan agen untuk melakukan penyaringan (\textit{filtering}) berbasis token kata terhadap ruang pencarian, sehingga hasil yang diperoleh lebih relevan dengan konteks tanaman atau penyakit yang sedang dianalisis. Skema lengkap representasi data pada indeks basis pengetahuan disajikan pada Tabel \ref{tab:knowledge_base_schema}.

\begin{table}[H]
    \centering
    \caption{Skema Poin dengan Vektor Jamak dan Metadata pada Basis Data Pengetahuan}
    \label{tab:knowledge_base_schema}
    \small
    \setstretch{1.0}
    \resizebox{\textwidth}{!}{%
    \begin{tabular}{|l|l|l|p{7cm}|}
    \hline
    \textbf{Komponen} & \textbf{Kunci} & \textbf{Tipe} & \textbf{Deskripsi} \\ \hline
    \textbf{Vektor Padat} & \texttt{gemini} & 3.072 & Representasi semantik dari Gemini Embedding. \\ \hline
    & \texttt{openai} & 3.072 & Representasi semantik dari OpenAI Embedding. \\ \hline
    \textbf{Vektor Jarang} & \texttt{splade} & Variabel & Representasi leksikal dari SPLADE++. \\ \hline
    & \texttt{bm25} & Variabel & Skor relevansi berbasis frekuensi kata (BM25). \\ \hline
    \textbf{Payload} & \texttt{page\_content} & String & Konten teks lengkap dari dokumen. \\ \hline
    & \texttt{metadata.plant} & String & Nama tanaman. Diindeks secara \textit{full-text}. \\ \hline
    & \texttt{metadata.disease} & String & Nama penyakit. Diindeks secara \textit{full-text}. \\ \hline
    & \texttt{metadata.type} & String & Jenis dokumen (budidaya/penyakit). \\ \hline
    & \texttt{metadata.source} & String & Sumber dokumen (PlantVillage/Gardenology). \\ \hline
    \end{tabular}%
    }
\end{table}
%TODO(Human): Perlu di-update visualisasi ini yang ngeliatin semua mode pencariannya (image-text, image-image, text-text, text-image)

\paragraph{Strategi Pencarian}
Implementasi strategi pencarian pada modul ini dirancang untuk memaksimalkan relevansi dokumen melalui pendekatan bertingkat. Proses pencarian diawali dengan memberikan akses agen LLM untuk menerapkan filter metadata (seperti nama tanaman, penyakit, atau jenis dokumen) untuk membatasi ruang pencarian hanya pada entitas yang relevan. Selanjutnya, sistem melakukan pencarian hibrida yang menggabungkan hasil dari representasi vektor padat (\textit{dense embedding}) dan vektor jarang (\textit{sparse embedding}).

Pada tahap pencarian hibrida, sistem dikonfigurasi agar agen LLM dapat melakukan pengambilan awal (\textit{prefetching}) terhadap sejumlah kandidat dokumen yang lebih besar, yakni sebesar $2k$, dari masing-masing indeks vektor. Hasil dari kedua jalur pencarian ini kemudian disatukan menggunakan algoritma \textit{Reciprocal Rank Fusion} (RRF). Algoritma ini memberikan bobot skor berdasarkan peringkat kemunculan dokumen di setiap daftar hasil, dengan tujuan memprioritaskan dokumen yang secara konsisten dinilai relevan oleh kedua metode pencarian. Persamaan \textit{Reciprocal Rank Fusion} (RRF) didefinisikan sebagai berikut:

\begin{equation}
    score(d \in D) = \sum_{r_d \in R(d)} \frac{1}{k + r_d}
\end{equation}

\noindent di mana $D$ adalah himpunan dokumen hasil, $R(d)$ adalah himpunan peringkat dokumen $d$, dan $k$ adalah konstanta peredam.

Dokumen yang telah digabungkan melalui RRF kemudian diproses lebih lanjut oleh komponen \textit{reranker}. Komponen ini mengevaluasi ulang relevansi semantik antara kueri dan dokumen kandidat untuk menghasilkan urutan akhir yang lebih akurat. Dari proses ini, hanya $k$ dokumen teratas yang diambil dan disajikan kepada agen sebagai konteks informasi. Visualisasi alur kerja strategi pencarian ini dapat dilihat pada Gambar \ref{fig:knowledge_base_search_workflow}.

\begin{figure}[H]
    \centering
    \includegraphics[width=1\textwidth]{bab3/placeholderalatpencarianbasispengetahuan.png}
    \caption{Visualisasi Alur Kerja Strategi Pencarian Hibrida pada Basis Data Pengetahuan}
    \label{fig:knowledge_base_search_workflow}
\end{figure}

\paragraph{Evaluasi}
Evaluasi kinerja modul pencarian informasi dilakukan secara \textit{offline} menggunakan pustaka \texttt{ranx} oleh \textcite{bassaniRanxBlazingFastPython2022}. Evaluasi ini bertujuan untuk mengukur efektivitas berbagai strategi pencarian dan konfigurasi model yang telah diimplementasikan dalam menemukan dokumen relevan dari korpus basis data pengetahuan.

Persiapan evaluasi dimulai dengan penyusunan \textit{Query Relevance Judgments} (Qrels). Qrels merupakan standar kebenaran (\textit{ground truth}) yang memetakan setiap kueri dengan dokumen-dokumen yang dianggap relevan. Dalam penelitian ini, struktur Qrels direpresentasikan dalam format \textit{dictionary} Python sesuai dengan spesifikasi \texttt{ranx}, di mana setiap ID kueri memiliki daftar ID dokumen relevan beserta skor relevansinya.

\begin{verbatim}
qrels_dict = {
    "q_1": {
        "d_1": 1,
        "d_2": 1,
    },
    # ...
}
\end{verbatim}

Skor relevansi yang digunakan bersifat biner, di mana nilai 1 diberikan untuk dokumen yang relevan dan secara implisit nilai 0 untuk dokumen lainnya. Hal ini didasarkan pada fakta bahwa dataset QA sintetis dibangun dengan memasangkan kueri secara spesifik pada potongan dokumen sumber tertentu, sehingga relevansi dapat ditentukan secara deterministik.

Proses evaluasi dilakukan dengan membandingkan hasil pencarian sistem (\textit{Run}) terhadap Qrels menggunakan serangkaian metrik standar dalam pencarian informasi. Mengingat penggunaan model \textit{embedding} dan LLM yang mutakhir (\textit{State-of-the-Art}), evaluasi difokuskan pada peringkat atas (Top-K) yang ketat, yaitu dengan nilai $k=5$. Pemilihan $k$ yang kecil ini bertujuan untuk memastikan sistem benar-benar menyajikan informasi yang paling relevan pada urutan teratas, sehingga meminimalkan beban kognitif pengguna dalam memilah informasi.

Berikut adalah metrik-metrik yang digunakan dalam evaluasi:

\begin{enumerate}[noitemsep]
    \item \textit{Mean Reciprocal Rank} (MRR) \\
    MRR mengukur rata-rata kebalikan dari peringkat dokumen relevan pertama yang ditemukan. Metrik ini sangat krusial untuk jenis pertanyaan faktual di mana pengguna hanya membutuhkan satu jawaban yang benar secepat mungkin. Nilai MRR dihitung dengan persamaan:
    \begin{equation}
        MRR = \frac{1}{|Q|} \sum_{i=1}^{|Q|} \frac{1}{rank_i}
    \end{equation}
    di mana $|Q|$ adalah jumlah kueri dan $rank_i$ adalah posisi peringkat dokumen relevan pertama untuk kueri ke-$i$.

    \item \textit{Normalized Discounted Cumulative Gain} (nDCG@k) \\
    nDCG@k mengevaluasi kualitas pemeringkatan dengan mempertimbangkan posisi dokumen relevan dalam daftar hasil. Dokumen relevan yang muncul di posisi lebih tinggi akan memberikan kontribusi skor yang lebih besar. Metrik ini dihitung dengan membagi \textit{Discounted Cumulative Gain} (DCG) dengan \textit{Ideal DCG} (IDCG).

    Nilai DCG pada peringkat ke-$k$ didefinisikan sebagai akumulasi skor relevansi yang didiskon secara logaritmik:
    \begin{equation}
        DCG@k = \sum_{i=1}^{k} \frac{rel_i}{\log_2(i+1)}
    \end{equation}

    Sedangkan IDCG adalah nilai DCG maksimal yang dapat dicapai jika dokumen-dokumen relevan diurutkan pada posisi ideal (paling atas):
    \begin{equation}
        IDCG@k = \sum_{i=1}^{|REL_k|} \frac{rel_i}{\log_2(i+1)}
    \end{equation}

    Sehingga, nilai nDCG@k diperoleh melalui persamaan:
    \begin{equation}
        nDCG@k = \frac{DCG@k}{IDCG@k}
    \end{equation}
    Metrik ini digunakan untuk memastikan bahwa algoritma pencarian tidak hanya menemukan dokumen yang benar, tetapi juga mampu mengurutkannya berdasarkan tingkat kepentingannya.

    \item \textit{Recall}@k \\
    Recall@k mengukur proporsi dokumen relevan yang berhasil ditemukan dalam $k$ hasil teratas dibandingkan dengan total dokumen relevan yang tersedia untuk kueri tersebut.
    \begin{equation}
        Recall@k = \frac{|\{ \text{dokumen relevan} \} \cap \{ \text{dokumen terambil dalam top-}k \}|}{|\{ \text{total dokumen relevan} \}|}
    \end{equation}
    Metrik ini penting untuk mengevaluasi cakupan (\textit{coverage}) sistem dalam menyajikan berbagai perspektif informasi yang mungkin tersebar di beberapa dokumen.
\end{enumerate}

\subsubsection{Alat Pencarian Internet}
Alat pencarian internet diimplementasikan menggunakan Tavily Search API. Tavily merupakan mesin pencari yang dioptimasi khusus untuk sistem agentik (\textit{agentic systems}) dengan tujuan memberikan hasil pencarian yang efisien, cepat, dan konsisten \parencite{TavilyaiTavilypython2026}. Penggunaan API ini memungkinkan sistem untuk mengabstraksi serangkaian proses kompleks dalam pengambilan informasi daring, yang mencakup pencarian, \textit{web scraping}, penyaringan konten, pemeringkatan ulang (\textit{re-ranking}), hingga ekstraksi informasi yang paling relevan secara otomatis.

Konfigurasi parameter yang digunakan dalam inisialisasi alat pencarian ini dirinci pada Tabel \ref{tab:tavily_parameters}. Parameter \texttt{max\_results} diatur untuk mengembalikan lima dokumen teratas guna memberikan variasi sumber informasi yang cukup bagi agen. Topik pencarian diatur ke "general" untuk memastikan cakupan sumber yang luas. Selain itu, fitur pencarian gambar diaktifkan untuk memperkaya konteks visual yang dapat diberikan kepada pengguna.

\begin{table}[H]
    \centering
    \caption{Konfigurasi Parameter Tavily Search API}
    \label{tab:tavily_parameters}
    \small
    \setstretch{1.0}
    \begin{tabular}{|l|l|p{7cm}|}
    \hline
    \textbf{Parameter} & \textbf{Nilai} & \textbf{Deskripsi} \\ \hline
    \texttt{max\_results} & 5 & Jumlah maksimal hasil dokumen yang dikembalikan. \\ \hline
    \texttt{topic} & "general" & Kategori topik pencarian untuk cakupan sumber umum. \\ \hline
    \texttt{search\_depth} & "advanced" & Kedalaman pencarian untuk hasil yang lebih komprehensif. \\ \hline
    \texttt{include\_images} & \textit{True} & Mengembalikan tautan gambar yang relevan. \\ \hline
    \texttt{include\_image\_descriptions} & \textit{True} & Menyertakan deskripsi tekstual untuk gambar yang ditemukan. \\ \hline
    \end{tabular}
\end{table}

\subsection{Perancangan dan Pengembangan Antarmuka Sistem Agentik}
Antarmuka sistem dikembangkan sebagai media interaksi utama antara pengguna dan sistem agentik yang telah dibangun. Pengembangan ini bertujuan untuk memfasilitasi akses terhadap fitur identifikasi penyakit tanaman secara intuitif melalui platform berbasis web.

\subsubsection{Arsitektur Antarmuka}
Sistem antarmuka dibangun dengan memisahkan komponen sisi klien (\textit{frontend}) dan sisi server (\textit{backend}) guna memastikan skalabilitas dan kemudahan pemeliharaan. Arsitektur sistem secara keseluruhan terdiri dari dua komponen utama:

\begin{enumerate}[nosep]
    \item \textbf{Aplikasi Web (\textit{Frontend})}: Komponen ini merupakan aplikasi web berbasis Next.js yang mengadopsi \textit{Agent Chat UI} dari \textcite{LangchainaiAgentchatui2026}. Aplikasi ini bertanggung jawab untuk merender antarmuka percakapan, mengelola unggahan berkas, serta menampilkan visualisasi pemanggilan alat (\textit{tool calls}) secara waktu nyata (\textit{real-time}).
    \item \textbf{LangGraph Server (\textit{Backend})}: Komponen ini berfungsi sebagai infrastruktur yang menyerupai arsitektur tanpa peladen (\textit{serverless-like}) yang menjalankan logika agen LangGraph secara berstatus (\textit{stateful}). Server menyediakan titik akhir (API) yang dioptimalkan untuk komunikasi asinkron, manajemen utas (\textit{thread}), dan persistensi status percakapan.
\end{enumerate}

Seluruh komponen pendukung yang beroperasi secara lokal, seperti koleksi basis data vektor Qdrant serta model-model inferensi (YOLOv11, OWLv2, dan SCOLD), diintegrasikan dalam satu lingkungan kontainer menggunakan Docker. Integrasi ini bertujuan untuk menyederhanakan proses penerapan (\textit{deployment}) serta memastikan konsistensi operasional sistem pada berbagai lingkungan penampungan (\textit{hosting}). Visualisasi arsitektur sistem secara garis besar (\textit{high-level}) disajikan pada Gambar \ref{fig:high_level_architecture}.

\begin{figure}[H]
    \centering
    \begin{tikzpicture}[
        node distance=1.5cm,
        box/.style={rectangle, draw, align=center, minimum width=2.5cm, minimum height=1cm, font=\scriptsize},
        sub/.style={draw, dashed, inner sep=0.4cm}
    ]
    \node (ui) [box, fill=green!10] {Agent Chat UI\\(Next.js)};
    \node (lgs) [box, right=2cm of ui, fill=blue!10] {LangGraph Server};

    \node (models) [sub, below=of lgs] {
        \begin{tabular}{c}
        \textbf{Modul \& Model} \\
        YOLOv11, OWLv2, SCOLD
        \end{tabular}
    };

    \node (db) [box, right=2cm of lgs, fill=red!10] {Qdrant DB};

    \node (ext) [sub, above=of lgs] {
        \begin{tabular}{c}
        \textbf{Layanan Eksternal} \\
        Gemini, Tavily, Jina, R2
        \end{tabular}
    };

    \draw [<->, thick] (ui) -- (lgs);
    \draw [<->, thick] (lgs) -- (db);
    \draw [<->, thick] (lgs) -- (models);
    \draw [<->, thick] (lgs) -- (ext);
    \end{tikzpicture}
    \caption{Arsitektur High-Level Sistem Agentik}
    \label{fig:high_level_architecture}
\end{figure}

\subsubsection{Use Case Diagram}
\begin{figure}[H]
    \centering
    \begin{tikzpicture}[
        actor/.style={circle, draw, minimum size=1cm},
        uc/.style={ellipse, draw, align=center, minimum width=3cm, font=\scriptsize},
        node distance=0.5cm
    ]
    \node (user) [actor] {User};
    \node (uc1) [uc, right=2cm of user, yshift=2cm] {Mengirim Prompt};
    \node (uc2) [uc, below=of uc1] {Unggah Citra};
    \node (uc3) [uc, below=of uc2] {Melihat Hasil};
    \node (uc4) [uc, below=of uc3] {Visualisasi Deteksi};
    \node (uc5) [uc, below=of uc4] {Kelola Riwayat};
    \node (uc6) [uc, below=of uc5] {Transparansi Agen};

    \draw (user) -- (uc1);
    \draw (user) -- (uc2);
    \draw (user) -- (uc3);
    \draw (user) -- (uc4);
    \draw (user) -- (uc5);
    \draw (user) -- (uc6);
    \end{tikzpicture}
    \caption{Use Case Diagram Antarmuka Sistem Agentik}
    \label{fig:use_case_diagram}
\end{figure}

\subsubsection{Sequence Diagram}
\begin{figure}[H]
    \centering
    \begin{tikzpicture}[
        node distance=1.5cm,
        font=\scriptsize,
        actor/.style={rectangle, draw, fill=gray!20, minimum width=2cm, minimum height=0.5cm},
        line/.style={draw, -stealth}
    ]
    % Lifelines
    \node (user) [actor] {Pengguna};
    \node (ui) [actor, right=2cm of user] {Agent Chat UI};
    \node (lgs) [actor, right=2cm of ui] {LangGraph Server};
    \node (llm) [actor, right=2cm of lgs] {Gemini 3 Flash};

    \draw [dashed] (user) -- ++(0,-7);
    \draw [dashed] (ui) -- ++(0,-7);
    \draw [dashed] (lgs) -- ++(0,-7);
    \draw [dashed] (llm) -- ++(0,-7);

    % Messages
    \draw [line] ([yshift=-1cm]user.south) -- node[above] {Unggah \& Tanya} ([yshift=-1cm]ui.south);
    \draw [line] ([yshift=-1.5cm]ui.south) -- node[above] {Kirim Payload} ([yshift=-1.5cm]lgs.south);
    \draw [line] ([yshift=-2cm]lgs.south) -- node[above] {Inisialisasi} ([yshift=-2cm]llm.south);

    % ReAct Loop box
    \draw [fill=blue!5, opacity=0.3] (ui.south -| lgs.west) ++(-0.5,-2.5) rectangle ++(6.5,-3);
    \node at ([yshift=-2.8cm, xshift=1cm]lgs.south) {\textbf{Siklus ReAct}};

    \draw [line] ([yshift=-3.5cm]llm.south) -- node[above] {Tool Call} ([yshift=-3.5cm]lgs.south);
    \draw [line, dashed] ([yshift=-4cm]lgs.south) -- node[above] {Observation} ([yshift=-4cm]llm.south);

    \draw [line] ([yshift=-5cm]llm.south) -- node[above] {Final Response} ([yshift=-5cm]lgs.south);
    \draw [line, dashed] ([yshift=-5.5cm]lgs.south) -- node[above] {Streaming} ([yshift=-5.5cm]ui.south);
    \draw [line, dashed] ([yshift=-6cm]ui.south) -- node[above] {Tampilkan} ([yshift=-6cm]user.south);

    \end{tikzpicture}
    \caption{Diagram Urutan Interaksi Identifikasi Penyakit Tanaman}
    \label{fig:sequence_diagram}
\end{figure}

\subsubsection{Activity Diagram}
%TODO(AI): Kasih kalimat pengantar tentang swimlane activity diagram. Jelasin tanggung jawabnya ada siapa ajah.

%TODO(AI): Masukin total ada 3 activity diagram


\subsubsection{Pengujian Antarmuka (Black Box Testing)}
Hasil pengujian dirangkum dalam Tabel \ref{tab:black_box_testing}.

\begin{table}[H]
    \centering
    \caption{Hasil Pengujian Black Box Antarmuka Sistem Agentik}
    \label{tab:black_box_testing}
    \small
    \setstretch{1.0}
    \begin{tabular}{|l|l|p{4cm}|p{5cm}|}
    \hline
    ID & Fitur & Skenario & Hasil Diharapkan \\ \hline
    UC1 & Prompt & Ketik \& Kirim & Respons streaming. \\ \hline
    UC2 & Unggah & Paperclip / Drag & Preview \& Terkirim. \\ \hline
    UC3 & Hasil & Tunggu Inferensi & Diagnosis \& Saran muncul. \\ \hline
    UC4 & Viz & Klik Hasil Visual & Panel Artifact terbuka. \\ \hline
    UC5 & History & Pilih Sesi & Load riwayat \& state. \\ \hline
    UC6 & Tools & Klik Ikon Wrench & Tool calls terlihat. \\ \hline
    \end{tabular}
\end{table}

\subsection{Pengujian dan Evaluasi Sistem Agentik}
Penelitian ini mengadopsi pendekatan \textit{LLM-as-a-judge}. Pengujian dibagi menjadi empat eksperimen: VQA, OOD, Batasan Domain, dan Ablasi.

\subsubsection{Persiapan Data}
Dataset pengujian VQA disusun dari dataset PlantWild dengan total 267 entri (3 gambar per 89 kelas). Pasangan pertanyaan dan jawaban referensi dibangkitkan secara sintetis menggunakan Gemini 2.5 Flash.

\begin{enumerate}
    \item \textbf{Skenario 1: Gejala Samar} (Deskripsi non-teknis).
    \item \textbf{Skenario 2: Pertanyaan Langsung} (Hanya sebut spesies).
    \item \textbf{Skenario 3: Pertanyaan Umum} (Hanya tanya kondisi).
\end{enumerate}
